% In this file you should put all LaTeX macros and settings to be used both by
% the pdf version and the web version.
% This should be most of your macros.

% The theorem-like environments defined below are those that appear by default
% in the dependency graph. See the README of leanblueprint if you need help to
% customize this. 
% The configuration below use the theorem counter for all those environments
% (this is what the [theorem] arguments mean) and never resets it.
% If you want for instance to number them within chapters then you can add
% [chapter] at the end of the next line.
\newtheorem{theorem}{Theorem}
\newtheorem{proposition}[theorem]{Proposition}
\newtheorem{lemma}[theorem]{Lemma}
\newtheorem{corollary}[theorem]{Corollary}

\theoremstyle{definition}
\newtheorem{definition}[theorem]{Definition}

% Custom "axiom" environment for primitives we axiomatize rather than prove.
% We reuse the `proposition` counter so it shows up on the dep graph.
\newtheorem{axiom}[theorem]{Axiom}

% Math notation used throughout
\newcommand{\R}{\mathbb{R}}
\newcommand{\Vect}[1]{\mathbb{R}^{#1}}
\newcommand{\Mat}[2]{\mathbb{R}^{#1 \times #2}}
\newcommand{\pdiv}{\operatorname{pdiv}}
\newcommand{\pdivMat}{\operatorname{pdivMat}}
\newcommand{\HasVJP}{\mathsf{HasVJP}}
\newcommand{\HasVJPMat}{\mathsf{HasVJPMat}}
\newcommand{\HasVJPthree}{\mathsf{HasVJP3}}
% Foundation-flip notation (post the axiom-elimination work).
\newcommand{\fderiv}{\operatorname{fderiv}}
\newcommand{\basisVec}{\mathbf{e}}
\newcommand{\Differentiable}{\mathsf{Differentiable}}
\newcommand{\vjp}{\operatorname{vjp}}
\newcommand{\vjpMat}{\operatorname{vjpMat}}
\newcommand{\biPathMat}{\operatorname{biPathMat}}
\newcommand{\bnIstdBroadcast}{\operatorname{bnIstdBroadcast}}